\section{Related Work}
\label{sec:related-work}

% Original scaffold TODO: organize prior work by theme and explain differences.

\subsection{Cultural Knowledge and Evidence}
Cultural question answering tests knowledge that general visual benchmarks may not cover.
CulturalVQA contains 2,378 image--question pairs from 11 countries, with questions about food, clothing, drinks, rituals, and traditions \citep{nayak-etal-2024-benchmarking}.
CVQA expands multilingual coverage through locally authored visual questions and English versions \citep{romero2024cvqa}.
GIMMICK combines six tasks across three datasets to examine cultural knowledge through different modalities and geographic information \citep{schneider-etal-2025-gimmick}.
These benchmarks make cultural coverage an explicit evaluation concern.
Country or language indexes are useful for organizing data, but do not require every depicted element to have one exclusive national association.
Our proposed annotation instead attaches supported associations to individual cue instances.

Several works separate recognition from other forms of evidence.
Seeing Culture first tests visual cultural recognition and then artifact segmentation \citep{satar-etal-2025-seeing}.
BLEnD-Vis aligns region-to-entity and entity-to-region question formats across text and images \citep{tan-etal-2026-blend}.
MMAC evaluates cultural grounding with aligned text, image, and speech inputs, including consistency across languages and modalities \citep{zheng-etal-2026-mmac}.
Our cue-modality task identifies the channels in a particular video that support a specified cue.

\subsection{Cultural Video Understanding and Generation}
ViMUL-Bench evaluates multilingual video understanding with multiple-choice and open-ended questions across cultural and general categories \citep{shafique-etal-2025-culturally}.
MINERVA-Cultural, published as CURVE, focuses on long-video reasoning with human-authored native-language questions, answers, and reasoning traces across 18 locales \citep{Singh_2026_CVPR}.
% arXiv:2601.10649v2 uses MINERVA-Cultural; the verified CVPR entry uses CURVE.
% Paper and repository QA counts differ; see notes/REFERENCE_AUDIT_CULTURE.md.
These benchmarks motivate evaluation beyond English and short, visually obvious events.
Our proposed tasks complement their question-answering settings by explicitly separating cue discovery, evidence attribution, and relations between elements.

The Cultural Moment Benchmark separates concept naming, visual recognition, and temporal localization \citep{satar-etal-2026-cultural}.
Its stages select a concept name from four options, select a matching video moment, and predict start and end times in a different example video.
Its 306 concepts cover seven Southeast Asian countries.
Our candidate grounding task allows multiple intervals for a recurring cue.

CultureVidBench evaluates text-to-video generation using 1,000 prompts spanning 12 countries and 14 cultural aspects \citep{han-etal-2026-culturevidbench}.
Its evaluation considers visual, textual, and audio cultural details.
When Cultures Move studies mono-cultural and cross-cultural generated videos, including combinations of people, actions, and locations \citep{li-etal-2026-cultures}.
\benchmark{} focuses on how models interpret cultural cues and their relations in collected videos.

\subsection{Cultural Mixing and Compositional Evaluation}
CultureMix directly studies cultural mixing in images \citep{Kim_2026_CVPR}.
It contains approximately 23,000 generated and human-verified food images in four settings: food alone, multiple foods, food with a cultural background, and multiple foods with a cultural background.
Its evaluation tests whether individual food and cultural labels remain consistent across these settings.
This is the closest motivation for preserving component identity when several cultural cues occur together.
CulMix-Video extends the research question to temporal and audiovisual evidence, including changes in social setting or carrier without requiring technical fusion.
Its candidate role task additionally asks how the elements participate in the observed event.

Winoground tests whether models distinguish image--caption pairings when captions use the same words in different orders \citep{thrush_and_ross2022winoground}.
This motivates role classification controls that keep the elements similar while varying their relation.

\subsection{Visual Relations and Semantic Roles}
Visual relationship detection predicts object--predicate--object relations in images \citep{lu2016visual}.
Visual Genome provides dense object, attribute, and relation annotations \citep{krishna2017visualgenome}.
ImageNet-VidVRD and VidOR extend relation annotation to videos with object trajectories and temporal intervals \citep{shang2017video,shang2019annotating}.

Semantic role labeling offers a complementary model.
imSitu represents an activity through role-filled participants \citep{Yatskar_2016_CVPR}. SWiG adds spatial grounding \citep{pratt2020grounded}.
VidSitu extends semantic roles to video, including entity coreference and relations between events \citep{Sadhu_2021_CVPR}.
M2E2 studies event and argument extraction across news images and text \citep{li-etal-2020-cross}.
For CulMix-Video, a component and its base artifact, or a practice and its carrier, occupy different roles.
